\documentclass[10pt]{article}
\providecommand{\FIGDIR}{fig}
\newcommand{\DIGESTDATE}{7 September 2026}
% ===================== PM Research Watch - shared preamble =====================
\usepackage[T1]{fontenc}
\usepackage[utf8]{inputenc}
\usepackage{amsmath,amssymb}
\usepackage{tgpagella}
\usepackage{tgheros}
\renewcommand{\sfdefault}{qhv}
\renewcommand{\ttdefault}{lmtt}
\usepackage{microtype}
\usepackage[english]{babel}

\usepackage[a4paper,top=17mm,bottom=17mm,left=17mm,right=17mm,
            headsep=6mm,footskip=10mm]{geometry}
\usepackage{graphicx}
\usepackage{xcolor}
\usepackage{colortbl}
\usepackage{tikz}
\usetikzlibrary{calc}
\usepackage[most]{tcolorbox}
\usepackage{titlesec}
\usepackage{fancyhdr}
\usepackage{booktabs}
\usepackage{array}
\usepackage{longtable}
\usepackage{enumitem}
\usepackage{multicol}
\usepackage{ragged2e}
\usepackage[hidelinks]{hyperref}

% ------------------------------------------------------------------ palette
\definecolor{Ink}    {HTML}{14262E}
\definecolor{Deep}   {HTML}{0F4C5C}
\definecolor{Teal}   {HTML}{1F7A8C}
\definecolor{Sky}    {HTML}{4FA3C4}
\definecolor{Sage}   {HTML}{6FA287}
\definecolor{Amber}  {HTML}{E3A008}
\definecolor{Clay}   {HTML}{B5643C}
\definecolor{Coral}  {HTML}{D1495B}
\definecolor{Violet} {HTML}{6A5B9E}
\definecolor{Slate}  {HTML}{4A6572}
\definecolor{Paper}  {HTML}{FBF9F5}
\definecolor{Mist}   {HTML}{EEEAE2}
\definecolor{Rule}   {HTML}{D9D3C9}

\pagecolor{Paper}
\color{Ink}

\hypersetup{colorlinks=true, linkcolor=Deep, urlcolor=Teal, citecolor=Teal,
            pdftitle={PM Research Watch}, pdfauthor={Automated daily digest}}

% ------------------------------------------------------------------ headings
\titleformat{\section}
  {\sffamily\bfseries\large\color{Deep}}{}{0pt}
  {}[{\vspace{1.2mm}\color{Amber}\titlerule[1.1pt]}]
\titlespacing*{\section}{0pt}{5.5mm}{2.8mm}

\titleformat{\subsection}
  {\sffamily\bfseries\normalsize\color{Teal}}{}{0pt}{}
\titlespacing*{\subsection}{0pt}{3.6mm}{1.4mm}

\setlist[itemize]{leftmargin=4.2mm, itemsep=1.1mm, topsep=1.3mm, parsep=0pt,
                  label=\textcolor{Teal}{\footnotesize$\blacktriangleright$}}

% ------------------------------------------------------------------ header/footer
\pagestyle{fancy}
\fancyhf{}
\renewcommand{\headrulewidth}{0pt}
\renewcommand{\footrulewidth}{0pt}
\fancyhead[L]{\sffamily\scriptsize\color{Slate}PM RESEARCH WATCH\;\textcolor{Rule}{\textbar}\;\DIGESTDATE}
\fancyhead[R]{\sffamily\scriptsize\color{Slate}Particulate matter monitoring \& health effects}
\fancyfoot[C]{\sffamily\scriptsize\color{Slate}\thepage}

% ------------------------------------------------------------------ components
% metric tile
\newcommand{\metric}[3]{%
  \begin{tcolorbox}[enhanced, colback=white, colframe=Rule, boxrule=0.5pt,
    arc=1.6mm, left=2.2mm, right=2.2mm, top=1.8mm, bottom=1.8mm,
    borderline west={1.7mm}{0pt}{#3}, nobeforeafter, valign=center,
    height=17mm, width=\linewidth]
    \RaggedRight
    {\sffamily\bfseries\fontsize{17}{17}\selectfont\color{#3}#1}\\[0.5mm]
    {\sffamily\fontsize{6.5}{7.4}\selectfont\color{Slate}#2}
  \end{tcolorbox}}

% pull-quote / key insight box
\newtcolorbox{keybox}[1][]{enhanced, colback=Mist, colframe=Mist,
  arc=1.6mm, left=3.2mm, right=3.2mm, top=2.4mm, bottom=2.4mm,
  borderline west={1.7mm}{0pt}{Amber}, #1}

% subtopic band
\newcommand{\band}[2]{%
  \vspace{3.2mm}\noindent
  \begin{tcolorbox}[enhanced, colback=#2, colframe=#2, arc=1.2mm,
    left=2.6mm, right=2.6mm, top=1.1mm, bottom=1.1mm, boxrule=0pt,
    width=\linewidth, nobeforeafter]
    {\sffamily\bfseries\small\color{white}#1}
  \end{tcolorbox}\vspace{1.6mm}}

% one paper entry: {short cite}{title}{journal}{doi}{body}
\newcommand{\paperentry}[5]{%
  \par\vspace{1.9mm}\noindent
  {\sffamily\bfseries\footnotesize\color{Ink}#1}\;%
  {\sffamily\fontsize{7}{8}\selectfont\color{Slate}#3}\\[0.4mm]
  {\footnotesize\itshape\color{Deep}#2}\\[0.6mm]
  {\footnotesize #4}\\[0.3mm]
  {\fontsize{6.6}{7.6}\selectfont\color{Slate}\texttt{doi:}\,\href{https://doi.org/#5}{\texttt{\detokenize{#5}}}}
  \par}

% caption style
\newcommand{\figcap}[1]{\par\vspace{1.0mm}
  {\fontsize{7.2}{8.6}\selectfont\color{Slate}\textsf{\textbf{Fig.}}\; #1\par}\vspace{1.5mm}}

\providecommand{\permil}{\textperthousand}

\setlength{\parindent}{0pt}
\setlength{\parskip}{0pt}
\renewcommand{\arraystretch}{1.18}


\begin{document}

% ==================================================================== MASTHEAD
\thispagestyle{empty}
\begin{tikzpicture}[remember picture, overlay]
  \fill[Deep] (current page.north west) rectangle
      ($(current page.north east)+(0,-67mm)$);
  \fill[Amber] ($(current page.north west)+(0,-67mm)$) rectangle
      ($(current page.north east)+(0,-68.4mm)$);
  \foreach \x/\y/\r in {22/12/0.9, 41/26/1.5, 63/9/0.7, 88/31/1.1, 112/17/1.9,
                        138/29/0.8, 159/13/1.3, 178/54/1.0, 196/21/1.6,
                        30/44/0.6, 74/57/1.2, 125/48/0.9, 168/6/0.8, 191/59/0.7,
                        52/18/1.0, 100/60/0.6, 148/20/1.1, 66/49/1.4, 205/45/0.8}
    \fill[Sky, opacity=0.30] ($(current page.north west)+(\x mm,-\y mm)$) circle (\r mm);
\end{tikzpicture}

\vspace*{4mm}
{\sffamily\fontsize{7.6}{9}\selectfont\color{Sky}%
 AUTOMATED DAILY LITERATURE DIGEST\;\textbullet\;PhD RESEARCH WATCH\par}
\vspace{2.2mm}
{\sffamily\bfseries\fontsize{27}{29}\selectfont\color{white}%
 Particulate Matter\par}
{\sffamily\fontsize{27}{29}\selectfont\color{Sky}%
 Monitoring \& Health Effects\par}
\vspace{5.0mm}
{\sffamily\fontsize{8.4}{10}\selectfont\color{white}%
 Issue 2026-09-07\;\;\textcolor{Amber}{\textbar}\;\;Entry window 7 September 2026\;\;%
 \textcolor{Amber}{\textbar}\;\;20 records screened in scope\par}
\vspace{18mm}

% -------------------------------------------------------------- metric strip
\noindent
\begin{minipage}[t]{0.190\linewidth}\metric{20}{RECORDS IN SCOPE\\(17 EXCLUDED)}{Deep}\end{minipage}\hfill
\begin{minipage}[t]{0.190\linewidth}\metric{8}{EFFECT ESTIMATES\\EXTRACTED}{Teal}\end{minipage}\hfill
\begin{minipage}[t]{0.190\linewidth}\metric{8}{WITH A HUMAN\\HEALTH ENDPOINT}{Sky}\end{minipage}\hfill
\begin{minipage}[t]{0.190\linewidth}\metric{8}{ON THE SENSING /\\EXPOSURE AXIS}{Amber}\end{minipage}\hfill
\begin{minipage}[t]{0.190\linewidth}\metric{7}{FROM JOURNALS\\PUBMED OMITS}{Clay}\end{minipage}

\vspace{6mm}

% -------------------------------------------------------------- signal of day
\section*{\sffamily\bfseries\color{Deep}Signal of the day}
\vspace{-1.5mm}{\color{Amber}\rule{\linewidth}{1.1pt}}\vspace{2.5mm}

\begin{keybox}
{\sffamily\bfseries\small\color{Deep}A regulatory station and a mobile monitor 250\,m apart
disagree enough to publish a different AQI colour one time in ten --- and distance does not
predict when.}\\[1.4mm]
\footnotesize
Paired fixed-station and simultaneous mobile PM$_{2.5}$ measurement at three contrasting
sites across the Boston metro area, spring and summer 2025 (\textit{Environ Sci Technol}),
puts a number on what a sparse network cannot see. Typical agreement is good --- mean
station-to-mobile difference \textbf{2.0\,$\mu$g/m$^3$} (median 1.5, IQR 0.6--2.8) --- but
the tail is not: dynamic events produced discrepancies exceeding \textbf{300\% in relative
terms and 800\,$\mu$g/m$^3$ in absolute terms}, and in \textbf{10.4\% of cases}, rising to
\textbf{21.8\% at one site in spring}, the two instruments implied \emph{different AQI
categories}. The finding that matters operationally is the negative one: \textbf{Euclidean
distance alone could not explain the differences}. Spatial representativeness is therefore
not a radius, which is exactly the assumption embedded in every siting guideline and every
inverse-distance interpolation built on top of a reference network. Relative humidity was
the strongest single correlate the abstract reports ($r=-0.84$) --- note the sign is against
a signed centre-minus-mobile difference, not an absolute one, so it is a direction, not a
magnitude, and the humidity result would need the full text to interpret. Two seasons, one
metro area, and mobile measurement carries its own inlet and response-time artefacts. Even
so, this is the cleanest published quantification of the gap that dense low-cost deployment
exists to close, and it converts that case from a spatial-coverage argument into a
public-messaging one.
\end{keybox}

\vspace{3mm}

% -------------------------------------------------------------- exec brief
\section*{\sffamily\bfseries\color{Deep}What changed today}
\vspace{-1.5mm}{\color{Amber}\rule{\linewidth}{1.1pt}}\vspace{2.5mm}

\footnotesize
\begin{itemize}
\item \textbf{The sensing-axis fix worked, and it worked completely.} Yesterday's action item
was to AND a particle clause into the connector's sensing query, which had returned 13 hits
of which 10 were rejected as topically irrelevant. Rewritten, it returned \textbf{10 hits,
all 10 particle-bearing, with zero topical false positives}. The phycotoxin, football-heading
and evapotranspiration class of error is gone. Recall did not suffer: the axis still
surfaced every sensing record the health axis missed.

\item \textbf{Crossref-by-ISSN carried this issue --- 7 of 20 records, from journals PubMed
structurally cannot index.} Yesterday the same leg returned 11 raw and \textbf{0} in scope.
Today it returned 46 raw and \textbf{7} in scope, including the signal of the day and three
\textit{Atmos Chem Phys} papers. This is the largest contribution this leg has made, and it
is a direct vindication of keeping it: \textit{ACP} and \textit{AMT} are invisible to the
PubMed spine, so a PubMed-only watch would have shipped 13 records and missed the one that
matters most to this project.

\item \textbf{Eleven of twenty records carry no human health endpoint --- the inverse of
yesterday.} 6 September was 10 of 15 \emph{with} an endpoint. Today's issue is dominated by
instrumentation and atmospheric process work: a UAV single-particle sampler, a tethered
balloon, two dust emission schemes, a forecasting network and a bioaerosol sampler
validation. Read this as a composition swing driven by which source leg fired, not as a
shift in the field.

\item \textbf{The trial-watch zero is a US federal holiday, not a weekend.} ClinicalTrials.gov
returned \textbf{0 studies updated registry-wide} on 7 September --- confirmed with a no-term
control. 7 September 2026 is \textbf{Labor Day}. The prior reading was that registry zeros are
a weekend artefact; this extends it to federal holidays and rules out the weekend framing as
the complete explanation. \texttt{trials.json} is unchanged at 11 entries.

\item \textbf{The forest plot is thin and structurally narrow: 8 estimates, 2 papers.} Three
are myopia prevalence ratios, five are AECOPD incidence rate ratios, and three of those five
are not particulate at all (ozone twice, diurnal temperature range once). The panel shows
within-study precision, not between-study agreement, and is captioned as such.

\item \textbf{Sulfate, not mass, carries the myopia signal in 212,566 children.} Tianjin
grades 4--6, machine-learning exposure surfaces at school locations: SO$_4^{2-}$ top vs
bottom quartile gives \textbf{PR 1.10 (1.07--1.13)} while the whole-mixture qgcomp estimate
is only \textbf{1.017 (1.007--1.027)} per quartile. A component-resolved contrast that
survives when unspeciated mass barely moves is the recurring argument against mass-only
networks --- and school-location assignment, not residence, is the design's own main threat.

\item \textbf{Three pipeline defects closed in the same run.} The DOI gate gained an
\texttt{UNREGISTERED-DOI} warn class after an \textit{Environ Res} record FAILed with a DOI
that PubMed asserts but Elsevier has not yet deposited (Crossref \emph{and} the Handle API
both 404). \texttt{plots.py} gained 17 new design labels --- the entire issue had fallen
through to \textit{Other / mixed} --- and the geography map now resolves the labels that were
silently pooling into the \textit{Global / multi-region} fallback, which had swallowed the
only \textbf{Sub-Saharan Africa} and only \textbf{Latin America} records in the window.
\end{itemize}

% ==================================================================== ANALYTICS
\section*{\sffamily\bfseries\color{Deep}Corpus analytics}
\vspace{-1.5mm}{\color{Amber}\rule{\linewidth}{1.1pt}}\vspace{3mm}

\noindent\includegraphics[width=\linewidth]{\FIGDIR/f1_subtopics.png}
\figcap{Subtopic assignment of the 20 in-scope records. Eight of the ten clusters are
occupied; \textit{Cardiovascular \& metabolic} and \textit{Neuro / mental health} are empty
for the first time in this archive's recent history. \textit{Exposure assessment} and
\textit{Sensing, forecasting \& instrumentation} tie at four each and together take 8 of 20
--- the methodological half of the watch outweighing the epidemiological half, which is
unusual and is a source-leg effect.}

\vspace{2mm}
\noindent\includegraphics[width=\linewidth]{\FIGDIR/f2_design_endpoint.png}
\figcap{Left --- study architecture. \textbf{Measurement campaigns dominate at 7 of 20},
against 5 for observational designs of all kinds combined. The campaign bucket is broad by
construction: it holds a UAV deployment, a tethered balloon, a citizen-science snow and dust
survey, a mobile-versus-fixed co-measurement, a six-city indoor survey, an event-based
bioaerosol watch and a hybrid low-cost/BAM network. None of these is a sensor co-location
study, which is a distinct and stricter claim. Right --- health endpoint. \textbf{Eleven of
20 carry no health endpoint at all}; of the nine that do, one is in-vitro toxicology, so the
human-endpoint count is eight.}

\vspace{2mm}
\noindent\includegraphics[width=\linewidth]{\FIGDIR/f5_forest.png}
\figcap{All eight quantitative estimates in the corpus, on a common logarithmic axis.
\textbf{They come from only two papers and the contrasts are not mutually comparable} ---
three are prevalence ratios for myopia against long-term quartile contrasts of a single
component, of a mixture and of ozone, and five are incidence rate ratios for COPD
exacerbation against 7-day mean lag0--6 windows. Three of the eight are not particulate:
ozone in summer and autumn, and a per-1$^\circ$C diurnal temperature range. Read direction
and precision only; the visual clustering here is within-study, not corroboration.}

\vspace{2mm}
\noindent\includegraphics[width=\linewidth]{\FIGDIR/f6_lifecourse.png}
\figcap{Life-course windows addressed; counts are non-exclusive and only 8 of the 20 records
place anywhere on this axis, because eleven have no human endpoint. The prenatal window is
occupied by both of today's cohort records; the adolescent window is empty again, as it has
been for most of this archive.}

\vspace{2mm}
\noindent\includegraphics[width=\linewidth]{\FIGDIR/f3_metric_geography.png}
\figcap{Left --- particle metric. \textit{PM (unspeciated) / emissions} is the largest bin at
\textbf{12 of 20}, and here that label is doing more work than usual: it holds modelled dust
emission flux, eBC and UV-absorbing PM, single-particle composition classes, PAH mass,
workplace metal aerosol and a virus-laden collection liquid. Only \textbf{8 of 20 report a
size-resolved PM metric}. Right --- geographic provenance. China supplies 6, and
\textbf{Sub-Saharan Africa and Latin America appear at one record each --- both were pooled
into the fallback bucket until the geography map was corrected in this run.}}

\vspace{2mm}
\noindent\includegraphics[width=\linewidth]{\FIGDIR/f4_heatmap.png}
\figcap{Subtopic $\times$ study-architecture cross-tabulation. \textbf{No cell holds more
than two records} --- four are tied at the maximum --- so this is a maximally dispersed day;
\textit{Exposure assessment} splits evenly between measurement campaigns and modelling.
The structural absence worth naming is
that no cell links an observational design to the sensing or exposure clusters --- no record
this window measures particles with a novel instrument \emph{and} attaches a health outcome
to what it measured. The Addis Ababa network is the nearest miss, and it reaches its health
claim through published risk coefficients rather than its own outcome data.}

% ==================================================================== DIGEST
\section*{\sffamily\bfseries\color{Deep}Paper-level digest}
\vspace{-1.5mm}{\color{Amber}\rule{\linewidth}{1.1pt}}\vspace{1mm}

{\fontsize{7.4}{9}\selectfont\color{Slate}Ordered by cluster. Each entry gives the
methodological core and the finding that matters, not an abstract paraphrase.\par}

\band{Sensing, forecasting \& instrumentation\hfill 4 records}{Amber}

\paperentry{Boston sparse network 2026}
{What sparse fixed air-quality networks miss, measured against mobile monitoring}
{Environ Sci Technol}
{Fixed-station PM$_{2.5}$ paired with simultaneous mobile measurement at a dense set of
locations within a 250\,m radius, at three contrasting areas of the Boston metro area across
spring and summer 2025. Typical agreement is close --- mean difference \textbf{2.0\,$\mu$g/m$^3$},
median 1.5, IQR 0.6--2.8 --- but dynamic events drove differences past \textbf{300\%} and
\textbf{800\,$\mu$g/m$^3$}, and \textbf{10.4\% of cases} (21.8\% at one site in spring) fell
either side of an AQI category boundary. Crucially \textbf{Euclidean distance alone could not
explain the discrepancies}, which undercuts radius-based representativeness and every
inverse-distance product built on a reference network. Relative humidity was the strongest
reported correlate ($r=-0.84$), a signed rather than absolute association that the abstract
does not fully disambiguate. Two seasons and one metro area, and mobile platforms carry inlet
and response-time artefacts of their own; but this is the number the dense-network argument
has been missing.}
{10.1021/acs.est.6c03923}

\paperentry{UAV single-particle CCSEM 2026}
{Size-resolved aerosol composition and mixing state through the boundary layer from a UAV}
{Environ Sci Technol}
{A microscopy-based particle sampler flown on a UAV at up to 30\,m vertical and 1\,min
temporal resolution in a coastal boundary layer; over \textbf{100{,}000 particles} classified
by computer-controlled SEM with EDX. The observed transition runs from sea-breeze aerosol
dominated by aged sea salt and dust (\textbf{46--83\%}) to a continentally influenced
population dominated by organic-sulfur particles (\textbf{57--89\%}), and the mixing state
index falls from \textbf{53\% to 27\%} as those particles form --- i.e.\ toward external
mixing. Mixing state showed strong diurnal and size dependence but weak vertical structure.
Offline microscopy is not a real-time measurement and the sampling is necessarily
intermittent, so the ``temporal evolution'' is a sequence of snapshots; one coastal site, and
no companion mass measurement is described. The transferable point for network design is that
composition and mixing state vary on scales a fixed filter sampler averages away entirely.}
{10.1021/acs.est.6c03627}

\paperentry{Bogota BiLSTM Huber 2026}
{Asymmetric-loss BiLSTM for 24\,h PM$_{2.5}$ forecasting and exceedance detection}
{Atmosphere}
{Hourly records from the Puente Aranda industrial-corridor station, Bogot\'a, 2022--2025
($n=31{,}961$). The motivating statistic is sharp: \textbf{98.7\%} of days exceeding the WHO
2021 15\,$\mu$g/m$^3$ guideline contain at least one hour above 25, median peak-to-mean
\textbf{1.90}, so a model trained on mean error systematically misses the episodes that
matter. Three mechanisms are stacked --- an asymmetric Huber loss penalising underprediction
$\alpha$-fold, exceedance sample weighting, and Optuna-TPE over 100 trials. The honest result
is a reported \textbf{regression--detection tension}: no $\alpha$ minimised both global and
peak MAE. Optimisation chose $\alpha=2.80$, giving the best detection (AP \textbf{0.724},
95\% CI 0.582--0.821) and lowest peak MAE (10.64\,$\mu$g/m$^3$) at the cost of
\textbf{$R^2\approx0$} globally; $\alpha=1.5$ is the compromise (AP 0.718, RMSE 6.65,
$R^2=0.372$). One station, so nothing here is a spatial claim, and the wide bootstrap
intervals on AP overlap across all configurations --- the ranking is weaker than the point
estimates suggest.}
{10.3390/atmos17090871}

\paperentry{Coriolis sampler inactivation 2026}
{Sampler-induced loss of respiratory virus infectivity in spiked collection liquid}
{Indoor Air}
{Rhinovirus, SARS-CoV-2 and influenza A spiked into the collection liquid of a Coriolis
$\mu$ liquid impinger and run at several airflow rates --- deliberately \emph{not} an aerosol
or ambient-air collection, which isolates mechanical stress in the liquid from capture
efficiency. Culturable virus fell by up to \textbf{98.6\%} for SARS-CoV-2 and influenza A,
while RNA copies were largely unchanged for SARS-CoV-2 and rhinovirus and reduced only for
influenza A. Rhinovirus infectivity was essentially unaffected, so the effect is
virus-specific rather than a general artefact. The design's strength is also its limit: it
cannot say how much infectivity is lost during aerosol capture itself, only during the liquid
phase, so it bounds one term of the total sampling loss. The consequence for bioaerosol
surveillance is direct --- a PCR-positive, culture-negative air sample is not evidence of
non-infectious aerosol.}
{10.1155/ina/5068026}

\band{Exposure assessment \& modelling\hfill 4 records}{Amber}

\paperentry{Liupan tethered eBC 2026}
{Vertical structure of carbonaceous aerosol over the Loess Plateau from tethered balloon}
{Atmos Chem Phys}
{Coordinated tethered-balloon profiling of carbonaceous aerosol and meteorology at Pingliang,
eastern foothills of the Liupan Mountains, during July 2023 and July 2024 campaigns.
Near-surface means were \textbf{0.84\,$\mu$g/m$^3$ eBC} and \textbf{1.24\,$\mu$g/m$^3$ UVPM},
both decreasing with height. The substantive observation is that in the early morning and at
night, when convection is weak, \textbf{UVPM aloft exceeded eBC} --- attributed to nucleation
from gaseous precursors rather than transport. Thermodynamic processes dominated vertical
transport near the surface while enhanced dynamics aloft drove horizontal dispersion, and
southerly air masses contributed disproportionately to UVPM. Two July campaigns at one site
cannot separate seasonal from interannual variation, and eBC and UVPM are optical
attributions from a filter-based absorption instrument, not chemical speciation --- the
nucleation interpretation is inference from the eBC/UVPM divergence, not a measured
formation pathway.}
{10.5194/acp-26-12671-2026}

\paperentry{EROD-HR/SOILHD Iberia 2026}
{High-resolution soil erodibility and its effect on WRF-Chem dust and PM$_{10}$}
{Atmos Chem Phys}
{Two erodibility datasets tested in WRF-Chem with GOCART: EROD-HR, which folds fine-resolution
topography in at 0.0625$^\circ$ globally and 1\,km over Iberia, and SOILHD, which adds
high-resolution sand/silt/clay texture and removes misclassified bare-soil areas at 1\,km
globally. Five dust episodes spanning 2022--2025 cover both local emission and long-range
North African intrusion, evaluated against \textbf{PM$_{10}$ from the SINQLAIR network in
Murcia}. Both datasets improved the magnitude and timing of modelled PM$_{10}$ peaks at
inland stations; \textbf{improvements were limited at coastal stations}, which the paper does
not fully resolve and which is the most interesting part of the result --- coastal sites mix
dust with sea salt and are the ones regulatory exceedance attribution most often disputes.
Five episodes is a small evaluation set, and network PM$_{10}$ cannot separate dust from
co-occurring anthropogenic mass.}
{10.5194/acp-26-12641-2026}

\paperentry{Taiwan LRT bioaerosol 2026}
{Long-range transport effects on airborne fungal concentration and persistence}
{Environ Res}
{Daily wintertime fungal monitoring in northern Taiwan, comparing periods before, during and
after Asian dust storms and frontal pollution events. Long-range transport raised fungal
concentrations by \textbf{57--193\% during events} and \textbf{17--69\% in the post-event
window}, with elevation persisting \textbf{one to four days} after passage. Frontal pollution
enhanced concentrations more broadly than dust events did, and the assemblages differed by
regime: dust storms enriched Basidiospores and \textit{Cladosporium}, frontal events enriched
\textit{Aspergillus}/\textit{Penicillium} and Ascospores. The persistence result is the useful
one, because event-triggered sampling designs that stop at event end will systematically
under-count exposure. Weaknesses are the usual ones for this design: event periods are
defined meteorologically and confounded with the temperature, humidity and mixing changes
that accompany them, and no source apportionment separates transported spores from locally
resuspended ones.}
{10.1016/j.envres.2026.125640}

\paperentry{St Elias proglacial dust 2026}
{Sub-kilometre modelling of mineral dust emission from proglacial valleys}
{Atmos Chem Phys}
{A sub-km WRF-Chem implementation built to resolve dust emission from the proglacial valleys
of the St.\ Elias Mountains, Canada --- sources absent from current emission modelling and
unaccounted for in global inventories. Land-surface inputs were rebuilt with explicit
treatment of glaciofluvial deposit heterogeneity and inundation state; 5--19\,d simulations
across 2019--2022 were evaluated against in situ meteorology, dust measurements, camera
stations and surface Doppler LiDAR. Total emission is estimated at \textbf{1.0$\times$10$^4$
kg\,km$^{-2}$\,d$^{-1}$} from erodible deposits over 47 simulated days, offered explicitly as
a first-order figure. Skill in reproducing surface meteorology and in-valley vertical
dispersion is \textbf{season-dependent}, and the emission scheme remains sensitive to soil
texture, so the flux estimate inherits that sensitivity. High-latitude dust is a genuinely
under-measured PM source and this is a careful first quantification, but 47 days at one
mountain range does not support extrapolation.}
{10.5194/acp-26-12613-2026}

\band{Respiratory \& allergic\hfill 3 records}{Sage}

\paperentry{Shijiazhuang AECOPD 2026}
{Seasonal and short-term environmental associations with COPD exacerbation admissions}
{COPD}
{City-wide retrospective time-series of \textbf{22{,}468 AECOPD admissions}, 2020--2024, from
the municipal medical insurance database covering all tertiary, secondary and community
hospitals --- unusually complete ascertainment, though ED discharges without admission are
excluded. Mean age 71.7 (SD 9.1), male-to-female 3.7:1, winter peak 6{,}009 against summer
5{,}051. On the 7-day mean of lags 0--6, each 10\,$\mu$g/m$^3$ of winter PM$_{10}$ gave
\textbf{IRR 1.069 (1.052--1.085)} and NO$_2$ \textbf{1.089 (1.050--1.129)}; ozone was positive
in summer \textbf{1.035 (1.015--1.056)} and autumn \textbf{1.022 (1.005--1.039)}; a
1$^\circ$C wider diurnal range gave \textbf{1.055 (1.034--1.077)}. The authors are explicit
that the many lag $\times$ pollutant $\times$ season tests are exploratory and that the
prediction models are not a deployable early-warning system --- a welcome discipline, and the
reason the apparent protective multipollutant IRRs should not be read as findings. Collinear
pollutants in a single city cannot be separated by adjustment.}
{10.1080/15412555.2026.2723102}

\paperentry{Umea wood smoke chamber 2026}
{Controlled wood smoke exposure, sensory irritation and small airway mechanics}
{Environ Res}
{Randomised, double-blind crossover chamber study: \textbf{14} healthy non-smokers (11 M, 3 F,
mean age 26) exposed for 2\,h to wood smoke at \textbf{PM$_1$ 453 $\pm$ 13\,$\mu$g/m$^3$} and
to filtered air, three-week washout, with forced oscillometry (tremoFlo) alongside spirometry
and Borg symptom scoring. Wood smoke produced significant \textbf{nasal irritation
($p=0.003$)} with trends for cough and odour, but \textbf{no exposure-specific change} in any
spirometric or oscillometric resistance parameter (R5, R5--R19, R5--R20). Reactance and
resonant frequency moved with time irrespective of exposure, which is the control the design
exists to provide. With $n=14$ healthy young adults this is badly underpowered for small-airway
effects and the null should be read as such, not as evidence of no effect --- the authors say
as much. The methodological point stands independently: irritative symptoms preceded any
mechanical signal, so spirometry-only designs will under-detect early response.}
{10.1016/j.envres.2026.125629}

\paperentry{Disaster respiratory review 2026}
{Respiratory emergencies in disasters}
{Thorac Res Pract}
{Narrative review spanning earthquake silica and asbestos dust, wildfire fine particulate,
flood-related mould, volcanic ash and acidic gases, chemical release, blast barotrauma and
pandemic threats, synthesised from PubMed/MEDLINE, Web of Science and Scopus with EM-DAT, WHO
and UNDRR reports. It reports that acute respiratory infection is a leading cause of
disaster-related morbidity and mortality across age groups, that chronic respiratory patients
face elevated exacerbation risk, and that ARDS management remains hardest in
resource-limited settings. There is \textbf{no stated search protocol, no inclusion criteria
and no risk-of-bias appraisal}, so it indexes the literature rather than weighing it, and no
quantitative estimate here can be traced to a defined evidence base. Carried as a citation
anchor for the disaster-PM framing and for its explicit treatment of climate change as a
driver of rising burden --- not as evidence.}
{10.4274/ThoracResPract.2026.2026-6-13}

\band{Occupational \& indoor exposure\hfill 3 records}{Clay}

\paperentry{Europe indoor PAH children 2026}
{Indoor-air PAH levels and child cancer risk across six European cities}
{Environ Res}
{Twelve EPA PAHs in air and deposited dust from schools, homes and sports halls in Athens,
Barcelona, Colchester, Copenhagen, Helsinki and Lisbon --- four of the 16 target compounds
were dropped for analytical limitations, which is stated. Spatial variability is large:
medians above \textbf{50{,}000\,pg/m$^3$} where intensive residential wood burning and poor
ventilation coincide, against \textbf{under 10{,}000\,pg/m$^3$} in south-western Europe.
High-molecular-weight PAHs ran higher in schools than homes, consistent with occupancy and
accumulation of outdoor-derived particles, while low-molecular-weight species tracked
household practices. Estimating air concentrations from deposited dust \textbf{underestimated
HMW-PAHs}, which is a directly useful negative result for anyone tempted to substitute dust
for active sampling. Excess cancer risk over a five-year primary-school exposure period
spanned \textbf{10$^{-8}$ to 10$^{-5}$}, below the 10$^{-4}$ reference. Cross-sectional
sampling cannot capture seasonal wood-burning cycles.}
{10.1016/j.envres.2026.125643}

\paperentry{Rouyn-Noranda citizen dust 2026}
{Citizen-science measurement of smelter-derived metals in snow and indoor dust}
{Environ Monit Assess}
{Residents of \textbf{23 homes} 540--2{,}080\,m from the Horne Smelter --- Canada's only copper
smelter --- collected their own indoor dust and adjacent snow samples for Cu, Pb, Cd, As, Zn
and Cr. Accounting for distance \emph{and} wind direction, smelter location explained
\textbf{73\% of snow variance for Cu and Cr}, 58\% for Cd and Pb, 56\% for As and 44\% for Zn
($p<0.01$); for indoor dust the distance-only explanatory power was much lower at
\textbf{36\% (Cu) to 15\% (As)}. Copper served as a smelter tracer and correlated strongly
with the other metals in snow ($R^2>70\%$). Few residential characteristics predicted indoor
dust once distance was controlled. Twenty-three homes is a small sample, snow and settled
dust are deposition integrals rather than airborne concentrations, and participant-collected
samples carry uncontrolled handling variance. The value is the demonstration itself: a
participatory design produced a defensible source-attribution gradient in a community with a
documented exposure dispute.}
{10.1007/s10661-026-15900-2}

\paperentry{Bi/Bi2Te3 workplace ICP-OES 2026}
{ICP-OES determination of bismuth and bismuth telluride in workplace air}
{Zhonghua Lao Dong Wei Sheng Zhi Ye Bing Za Zhi}
{Method development for an occupational aerosol with no established Chinese workplace
standard: mixed cellulose ester filter sampling, ultra microwave digestion at 130$^\circ$C for
10 minutes, then simultaneous Bi and Te by ICP-OES with Bi$_2$Te$_3$ back-calculated from
either element using on-site speciation information. Linear over \textbf{0.04--30\,mg/L}
($r>0.9990$); on a 75\,L sample diluted to 50\,mL the minimum detectable Bi$_2$Te$_3$ is
\textbf{0.006\,mg/m$^3$} via Bi and \textbf{0.008} via Te. Digestion efficiency
99.1--100.0\% (Bi) and 101.2--103.3\% (Te), intra-batch precision under 0.8\%, 7-day room
temperature stability, and no interference from Ag, Al or As up to 10\,mg/L. Validated on
simulated samples rather than a field inter-laboratory comparison, and thermoelectric
workplaces are the only stated application. Included as the kind of low-visibility analytical
groundwork on which occupational exposure limits eventually rest.}
{10.3760/cma.j.cn121094-20250319-00107}

\band{Mechanistic toxicology\hfill 2 records}{Violet}

\paperentry{Fetal macrophage FCER1G 2026}
{Prenatal PM$_{2.5}$, fetal macrophage epigenetic reprogramming and atopic dermatitis}
{Nat Commun}
{Birth-cohort association between \textbf{first-trimester} PM$_{2.5}$ and early-childhood
atopic dermatitis, carried into mechanism by integrative multi-omics including single-cell
profiling of placenta and fetal immune cells. The central finding is \textbf{hypomethylation
of FCER1G in fetal macrophages} driving sustained overexpression that persists across
developmental stages and re-emerges in M2 macrophages of AD skin and peripheral blood;
functional work links FCER1G networks to NADPH-oxidase-mediated ROS and Th2 inflammatory
signalling. This is a rare case of a durable, cell-type-resolved epigenetic mediator proposed
for a prenatal PM$_{2.5}$ association rather than asserted. The exposure remains a modelled
residential estimate assigned to a trimester, which is the weakest link in an otherwise strong
chain --- first-trimester windows are precisely where exposure misclassification is largest.
The persistence claim rests on cross-sectional comparisons at different stages, not on
following the same cells.}
{10.1038/s41467-026-76298-6}

\paperentry{6PPD-Q AC16 p53 2026}
{Cardiotoxic mechanism of the tire-derived transformation product 6PPD-quinone}
{Ecotoxicol Environ Saf}
{Network toxicology and docking followed by wet-lab confirmation in human AC16
cardiomyocytes. \textbf{144 putative targets} with TP53 the top hub; \textbf{IC$_{50}$
13.92\,$\mu$M} at 24\,h. 6PPD-Q drove ROS accumulation and \textbf{p53 Ser15 phosphorylation
without raising total p53 protein or TP53 mRNA} --- i.e.\ post-translational stress
activation --- alongside TRAF6 upregulation, MAPK1 phosphorylation, IL6/TNF induction,
mitochondrial apoptosis (membrane potential loss, Bax/Bcl-2 rise, cleaved caspase-9 and -3)
and blocked late autophagic flux (LC3-II accumulation with p62 depletion). Pifithrin-$\alpha$
attenuated all of it, which is the load-bearing experiment. The exposure gap is total: a
micromolar dose to an immortalised cardiomyocyte line has no established relationship to any
airborne 6PPD-Q concentration, and no inhalation route is demonstrated. Relevant here only
because 6PPD-Q is a measurable constituent of road dust and traffic-derived PM.}
{10.1016/j.ecoenv.2026.120754}

\band{Burden, policy \& mitigation\hfill 2 records}{Coral}

\paperentry{Addis Ababa hybrid net 2026}
{PM$_{2.5}$ spatiotemporal variation and monetised burden from a hybrid sensor network}
{Health Serv Insights}
{Two years of hourly PM$_{2.5}$ (2022--2023) from a \textbf{hybrid network of seven low-cost
sensors and two beta attenuation monitors} --- the design this project cares about, deployed
in a city with almost no prior monitoring. Site means \textbf{15--33\,$\mu$g/m$^3$}, citywide
\textbf{27}, more than fivefold the WHO guideline of 5; concentrations were higher in the wet
Kiremt season, with a \textbf{07:00 diurnal peak of 38}. Hazard quotients exceeded 1 in all
age groups. The cancer-risk estimate is honestly reported as unstable: ELCR exceeded
10$^{-4}$ under a slope factor of 0.008 but stayed acceptable at 0.0012. Monetised DALYs give
\textbf{USD 652.95\,M} over two years, \textbf{0.25\% of GDP} in each year, dominated by
ischaemic heart disease. The paper does not describe co-location or humidity correction for
the seven low-cost nodes, and every health figure is transferred from external risk
coefficients, so the burden estimate inherits both instrument bias and transfer error.}
{10.1177/11786329261487535}

\paperentry{Hohhot park PM removal 2026}
{Random-forest attribution of PM reduction by urban park green space}
{Environ Pollut}
{Thirty-seven parks in Hohhot, summer daytime, reduction rates computed for \textbf{six PM
size fractions}. Linear, GAM and random-forest models were compared under leave-one-park-out
cross-validation and RF won, though the fit is modest and honestly reported --- point-level
$R^2$ \textbf{0.19--0.58}, park-level \textbf{0.18--0.72}. SHAP plus the derivatives of
partial-dependence and ICE curves identified nine drivers and, more usefully, the
\textbf{threshold ranges beyond which the marginal mitigation effect plateaus}: air
temperature 22.08--33.26$^\circ$C, RH 29.71--67.17\%, wind speed 0.14--1.12\,m/s, distance
from centre 2.7--16.79\,km, perimeter-area ratio 0.15--0.33, contiguity 0.24--0.48,
aggregation index 78.46--86.85\%. Summer daytime only, one semi-arid city, and the
in-versus-out contrast cannot separate deposition from dispersion and shelter effects. The
plateau structure is the transferable result; the specific thresholds are not.}
{10.1016/j.envpol.2026.129128}

\band{Other clinical endpoints\hfill 1 record}{Slate}

\paperentry{Tianjin myopia 2026}
{Long-term PM$_{2.5}$ components and childhood myopia}
{Environ Pollut}
{School-based study of \textbf{212{,}566} children in grades 4--6 in Tianjin, with three-year
mean PM$_{2.5}$, its major components (SO$_4^{2-}$, NO$_3^-$, NH$_4^+$, organic matter, black
carbon) and ozone estimated by machine-learning exposure models and assigned to \emph{school}
locations. Restricted cubic splines and quartile-based modified Poisson models for single
pollutants, quantile g-computation for the mixture. \textbf{Sulfate carried the signal}: top
versus bottom quartile \textbf{PR 1.10 (1.07--1.13)}, against a weak non-monotonic ozone
association \textbf{1.03 (1.00--1.05)} and a whole-mixture estimate of only \textbf{1.017
(1.007--1.027)} per quartile increase. Cross-sectional, so myopia onset timing is unknown and
reverse pathways through indoor time are not excluded; exposure at school ignores home and
commute, which for a three-year window is substantial misclassification. The component
contrast is the reason to read it --- a speciated signal that unspeciated mass does not
reproduce is the standing argument against mass-only networks.}
{10.1016/j.envpol.2026.129120}

\band{Reproductive \& developmental\hfill 1 record}{Sky}

\paperentry{SIP prenatal Bayley 2026}
{In utero PM$_{2.5}$ and NO$_2$ and cognitive-behavioural development at 1--3 years}
{Child Neuropsychol}
{Stress in Pregnancy cohort, \textbf{$n=210$}, Queens, New York; Bayley-III and ABAS-3 at 18,
24 and 36 months; prenatal PM$_{2.5}$ and NO$_2$ taken from public monitoring databases.
The reported pattern is \textbf{sex-specific and pollutant-specific}: PM$_{2.5}$ associated
with poorer communication, social skills and gross motor function in \textbf{boys}, NO$_2$
with poorer communication, gross motor and social-emotional skills in \textbf{girls}. No
effect estimates or intervals appear in the abstract, which is why nothing from this paper
enters the forest plot. With $n=210$ split by sex across three timepoints and multiple Bayley
and ABAS subdomains, the multiplicity burden is severe and the crossed sex-by-pollutant
pattern is exactly the shape that subgroup analysis produces by chance. Exposure from public
monitors rather than a spatial model adds classical misclassification. Hypothesis-generating,
and the authors largely frame it that way.}
{10.1080/09297049.2026.2723475}

% ==================================================================== TABLE
\section*{\sffamily\bfseries\color{Deep}Corpus register}
\vspace{-1.5mm}{\color{Amber}\rule{\linewidth}{1.1pt}}\vspace{2.5mm}

{\fontsize{6.7}{8.4}\selectfont
\setlength{\tabcolsep}{2.4pt}
\begin{longtable}{@{}>{\RaggedRight}p{25mm}>{\RaggedRight}p{41mm}>{\RaggedRight}p{27mm}%
>{\RaggedRight}p{23mm}>{\RaggedRight}p{18mm}>{\RaggedRight}p{15mm}@{}}
\toprule
\textsf{\textbf{First author}} & \textsf{\textbf{Focus}} & \textsf{\textbf{Journal}} &
\textsf{\textbf{Design}} & \textsf{\textbf{Metric}} & \textsf{\textbf{Region}}\\
\midrule
\endfirsthead
\toprule
\textsf{\textbf{First author}} & \textsf{\textbf{Focus}} & \textsf{\textbf{Journal}} &
\textsf{\textbf{Design}} & \textsf{\textbf{Metric}} & \textsf{\textbf{Region}}\\
\midrule
\endhead
\multicolumn{6}{@{}l@{}}{\cellcolor{Mist}\textsf{\textbf{\textcolor{Deep}{Neuro / mental health}}}}\\[0.6mm]
\href{https://doi.org/10.1001/jamanetworkopen.2026.25015}{Dales} & Wildfire-sourced PM$_{2.5}$ and ED visits for migraine and other primary headache & \textit{JAMA Netw Open} & Case-crossover & PM$_{2.5}$ (wildfire) & Canada\\
\href{https://doi.org/10.1093/ageing/afag215}{Glover} & Umbrella meta-review of urban environmental exposures and cognitive health & \textit{Age Ageing} & Umbrella review / meta-analysis & PM$_{2.5}$ & Global\\
\href{https://doi.org/10.1016/j.nsa.2026.107019}{Goossens} & In utero air pollution exposure and pre-/postnatal brain development & \textit{Neurosci Appl} & Narrative review & PM$_{2.5}$ / PM$_{10}$ / PM$_{\mathrm{coarse}}$ & Global\\
\href{https://doi.org/10.1016/j.isci.2026.116617}{Ma} & PM$_{2.5}$-induced neurotoxicity and depression: exposure risks and mechanisms & \textit{iScience} & Narrative review & PM$_{2.5}$ + components & Global\\
\href{https://doi.org/10.3760/cma.j.cn112150-20250804-00752}{Pan} & Maternal PM exposure and offspring neurodevelopment: effects and mechanisms & \textit{Zhonghua Yu Fang Yi Xue Za Zhi} & Narrative review & PM (bulk) & China\\
\addlinespace[0.9mm]
\multicolumn{6}{@{}l@{}}{\cellcolor{Mist}\textsf{\textbf{\textcolor{Teal}{Cardiovascular \& metabolic}}}}\\[0.6mm]
\href{https://doi.org/10.3760/cma.j.cn112338-20251120-00843}{Huang} & PM$_{2.5}$ and dyslipidaemia in rural Chinese children; dietary effect modification & \textit{Zhonghua Liu Xing Bing Xue Za Zhi} & Cross-sectional (multistage) & PM$_{2.5}$ & China\\
\href{https://doi.org/10.32394/pe/218966}{Kazmierska} & Air pollution and cardiovascular disease: review of current data & \textit{Przegl Epidemiol} & Narrative review & PM$_{2.5}$ / PM$_{10}$ & Global\\
\href{https://doi.org/10.1007/s00125-026-06814-2}{Xi} & Air pollution, built environment and temperature vs. prediabetes phenotypes & \textit{Diabetologia} & Prospective cohort & PM$_{2.5}$ & Germany\\
\href{https://doi.org/10.1161/JAHA.125.046917}{Yang} & Short-term PM$_{2.5}$ and hospitalisation/ED use in atrial fibrillation & \textit{J Am Heart Assoc} & Case-crossover & PM$_{2.5}$ (wildfire-influenced) & USA\\
\addlinespace[0.9mm]
\multicolumn{6}{@{}l@{}}{\cellcolor{Mist}\textsf{\textbf{\textcolor{Sky}{Reproductive \& developmental}}}}\\[0.6mm]
\href{https://doi.org/10.1016/j.jes.2025.10.043}{Chen} & Prenatal air pollution, fetal growth trajectory phenotypes and child neurodevelopment & \textit{J Environ Sci (China)} & Prospective cohort & PM$_{2.5}$ / PM$_{10}$ & China\\
\href{https://doi.org/10.1016/j.ecoenv.2026.120566}{Huang} & Joint prenatal bisphenol + air-pollutant mixture exposure and preschool IQ & \textit{Ecotoxicol Environ Saf} & Prospective cohort (mixtures) & PM$_{2.5}$ / PM$_{10}$ & China\\
\href{https://doi.org/10.1007/s40572-026-00552-8}{Peterson} & Ambient air pollution, PAHs and pubertal development: critical review & \textit{Curr Environ Health Rep} & Critical review & PM$_{2.5}$ components & Global\\
\href{https://doi.org/10.1016/j.envres.2026.125313}{Xie} & Long-term PM$_{2.5}$ and timing of pubertal development in boys and girls & \textit{Environ Res} & Prospective cohort & PM$_{2.5}$ & China\\
\addlinespace[0.9mm]
\multicolumn{6}{@{}l@{}}{\cellcolor{Mist}\textsf{\textbf{\textcolor{Sage}{Respiratory \& allergic}}}}\\[0.6mm]
\href{https://doi.org/10.1007/s00484-026-03289-z}{Bayar Muluk} & Climate change and otolaryngology: airway disease impacts and emission reduction & \textit{Int J Biometeorol} & Narrative review & PM$_{2.5}$ & Global\\
\href{https://doi.org/10.4103/mgr.MEDGASRES-D-25-00164}{Liao} & Short-term air pollution, respiratory symptoms and lung function in asthma & \textit{Med Gas Res} & Panel study & PM$_{2.5}$ / PM$_{10}$ & China\\
\href{https://doi.org/10.3390/epidemiologia7040097}{Sanchez de Prada} & Ambient PM and asthma-coded admissions in children 0-3 y, Spain & \textit{Epidemiologia} & Ecological panel & PM$_{2.5}$ / PM$_{10}$ & Spain\\
\href{https://doi.org/10.1016/j.jhazmat.2026.143076}{Yang} & Allergenic vs. toxic risk of airborne seed hairs compared with ambient PM$_{2.5}$ & \textit{J Hazard Mater} & Field measurement & PM$_{2.5}$ + bioaerosol & China\\
\addlinespace[0.9mm]
\multicolumn{6}{@{}l@{}}{\cellcolor{Mist}\textsf{\textbf{\textcolor{Violet}{Mechanistic toxicology}}}}\\[0.6mm]
\href{https://doi.org/10.1016/j.jes.2025.12.074}{Ding} & PM$_{2.5}$ promotes vascular calcification via miR-27b-5p/HES1 signalling & \textit{J Environ Sci (China)} & In vitro + organoid + murine & PM$_{2.5}$ & China\\
\href{https://doi.org/10.1007/s00360-026-01698-5}{Maraschi} & Osmoregulatory disruption and metal bioaccumulation from settleable atmospheric PM & \textit{J Comp Physiol B} & Controlled exposure (ecotox) & Settleable APM & Brazil\\
\href{https://doi.org/10.1016/j.jes.2025.11.018}{Wu} & Macrophage NLRP3 inflammasome in real-ambient PM-induced hepatic glucose dysregulation & \textit{J Environ Sci (China)} & Real-ambient murine + co-culture & Real-ambient PM & China\\
\addlinespace[0.9mm]
\multicolumn{6}{@{}l@{}}{\cellcolor{Mist}\textsf{\textbf{\textcolor{Clay}{Occupational \& indoor}}}}\\[0.6mm]
\href{https://doi.org/10.3389/fmedt.2026.1796828}{Chen} & Nested suction instrument to cut intraoperative surgical-smoke PM$_{2.5}$ & \textit{Front Med Technol} & Bench / device experiment & PM$_{2.5}$ (surgical smoke) & China\\
\href{https://doi.org/10.1093/annweh/wxag055}{Cole} & Determinants of respirable crystalline silica exposure in Australian tunnelling & \textit{Ann Work Expo Health} & Observational exposure study & Respirable dust / RCS & Australia\\
\href{https://doi.org/10.1016/j.jhazmat.2026.143081}{Ding} & Stage-resolved hair Zn-isotope dynamics as a firefighting combustion-exposure tracer & \textit{J Hazard Mater} & Repeated-measures biomarker & Combustion-derived PM & China\\
\href{https://doi.org/10.3390/jox16040119}{Ramsden} & Xenobiotic hazards in aircraft cabin air & \textit{J Xenobiot} & Review + risk economics & Ultrafine particles & Global\\
\addlinespace[0.9mm]
\multicolumn{6}{@{}l@{}}{\cellcolor{Mist}\textsf{\textbf{\textcolor{Amber}{Sensing, forecasting \& instrumentation}}}}\\[0.6mm]
\href{https://doi.org/10.1016/j.jenvman.2026.130515}{Shi} & Adaptive spatiotemporal graph transformer for multi-site multi-step PM$_{2.5}$ forecasting & \textit{J Environ Manage} & Deep-learning model & PM$_{2.5}$ & China / India\\
\href{https://doi.org/10.3390/bios16070380}{Tsuchiya} & Seasonal and indoor-outdoor characteristics of fluorescent aerosol particles in offices & \textit{Biosensors} & Real-time sensor campaign & Bioaerosol / coarse & Japan\\
\addlinespace[0.9mm]
\multicolumn{6}{@{}l@{}}{\cellcolor{Mist}\textsf{\textbf{\textcolor{Amber}{Exposure assessment \& modelling}}}}\\[0.6mm]
\href{https://doi.org/10.1029/2025GH001621}{Solorzano-Araque} & Multiscale exposure modelling: LOTOS-EUROS + mobility surveys + particle toxicity & \textit{GeoHealth} & CTM + mobility integration & PM$_{2.5}$ / PM$_{10}$ & Colombia\\
\addlinespace[0.9mm]
\multicolumn{6}{@{}l@{}}{\cellcolor{Mist}\textsf{\textbf{\textcolor{Coral}{Burden, policy \& mitigation}}}}\\[0.6mm]
\href{https://doi.org/10.1002/hsr2.72886}{Karamzad} & Ischaemic heart disease burden in MENA 1990-2021; PM attributable fraction 28.2\% & \textit{Health Sci Rep} & GBD secondary analysis & Ambient PM & MENA\\
\href{https://doi.org/10.1371/journal.pone.0347989}{Rees} & HEPA filtration to reduce respiratory infection in care homes & \textit{PLoS One} & RCT process evaluation & Filtered PM & UK\\
\href{https://doi.org/10.1016/j.jes.2026.01.009}{Shan} & Policy-driven emission reductions and post-pandemic rebound, high-resolution inventories & \textit{J Environ Sci (China)} & Emission inventory & PM emissions & China\\
\addlinespace[0.9mm]
\multicolumn{6}{@{}l@{}}{\cellcolor{Mist}\textsf{\textbf{\textcolor{Slate}{Other clinical endpoints}}}}\\[0.6mm]
\href{https://doi.org/10.3390/dj14070440}{Lupita} & PM$_{2.5}$, preventive-care capacity and caries/edentulism burden across the EU & \textit{Dent J} & Ecological panel & PM$_{2.5}$ & EU-27\\
\href{https://doi.org/10.1016/j.jdin.2026.06.008}{Park} & Long-term PM exposure and risk of androgenetic alopecia & \textit{JAAD Int} & Nationwide cohort & PM$_{2.5}$ / PM$_{10}$ & South Korea\\
\href{https://doi.org/10.3171/2026.3.PEDS2615}{Sun} & Environmental exposure and area deprivation vs. paediatric hydrocephalus outcomes & \textit{J Neurosurg Pediatr} & Registry cohort + geospatial & PM$_{2.5}$ (satellite-derived) & USA\\
\addlinespace[0.9mm]
\bottomrule
\end{longtable}}


\vspace{4mm}
\begin{keybox}
{\sffamily\bfseries\footnotesize\color{Deep}Provenance and method}\\[1.2mm]
{\fontsize{7.2}{8.8}\selectfont
Entry window \textbf{7 September 2026}, a single day, run against \texttt{[Date - Entry]}.
\textbf{PubMed} (NCBI E-utilities) returned \textbf{8 health-axis} and \textbf{6
sensing-axis} hits to the harvester; the \textbf{PubMed connector}, run independently on both
axes, returned \textbf{11} and \textbf{10} for 16 unique records. The sensing axis was
rewritten this run to AND a particle clause into the free-text terms, closing the action item
raised on 6 September: \textbf{all 10 hits were particle-bearing and none was rejected as
off-topic}, against 10 of 13 rejected yesterday. \textbf{Europe PMC}
(\texttt{CREATION\_DATE}) returned 9 raw and contributed 5 records neither PubMed leg had.
\textbf{Crossref-by-ISSN} returned \textbf{46 raw}, of which 13 were retrieved in full and
\textbf{7 shipped} --- the largest contribution this leg has made, and the source of the
signal of the day; \textit{Atmos Chem Phys}, \textit{Atmos Meas Tech} and \textit{Atmosphere}
are not indexed by the PubMed spine at all. \textbf{arXiv} returned a payload but nothing
inside the local recency screen. \textbf{OpenAlex} returned HTTP 429, unchanged since 28 July.
\textbf{Consensus} returned three hits, all far out of window (2019, 2020, 2025) --- a nil for
the \textbf{ninth} consecutive run. \textbf{Scholar Gateway} was queried with a narrowed
single-paper question after the 6 September budget overrun; it returned within budget but did
not hold the target article, so \texttt{10.1016/j.atmosenv.2026.122350} ships no summary and
joins the abstract retry list. \textbf{ClinicalTrials.gov} posted \textbf{no updates
registry-wide} on 7 September, verified with a no-term control; 7 September 2026 is
\textbf{Labor Day}, so the registry zero is a federal-holiday artefact and not evidence about
PM trial activity. \texttt{trials.json} is unchanged at 11.
Forty-four raw PubMed and Europe PMC hits deduplicated to \textbf{24 unique} on PMID $\to$
DOI $\to$ normalised title; \textbf{3 had already shipped in the 6 September issue} and were
dropped. Seventeen records were excluded with a logged reason --- those 3 duplicates, plus
aquatic and marine suspended particulate matter, gas-phase-only exposures, a CO$_2$ sensor, a
radiative-cooling materials study, tobacco-product chemistry and plant ecotoxicology.
The DOI gate passed 20 of 20 with two advisory warnings and gained a new
\texttt{UNREGISTERED-DOI} class this run: \texttt{10.1016/j.envres.2026.125643} is absent from
both Crossref and the Handle API but is asserted by PubMed's own record, so it is a publisher
deposit lag rather than a wrong identifier, and \textbf{that link may not resolve until
Elsevier deposits it}.
Subtopic, design, particle-metric and geography labels are assigned from title, abstract and
keywords; assignment is single-label by dominant contribution and is a judgement, not a
controlled vocabulary. Seventeen new design labels and several geography labels were added to
the plotting maps in this run. \textit{Global / multi-region} remains a fallback bucket that
absorbs unstated settings and bench work. Effect estimates are transcribed verbatim and use
heterogeneous contrasts --- see the forest-plot caption.
All DOI links resolve to the publisher of record. Abstracts remain the copyright of their
respective publishers.\par}
\end{keybox}

\end{document}
